%% ECON 282E -- Session 4, Deck A3: Optimization in Many Dimensions
%% Ported from QM Lec.4 F18-F33 (multivariate Newton, the curse, the BFGS worked
%% example) with the source's errors corrected -- see slides/SOURCES.md.
%% Numbers from tools/figures/s04_numbers.json, keys "newton_cost",
%% "bfgs_worked", "cg_worked", "derivative_free".

%% ECON 282E -- Foundations of Macroeconomics (UC Riverside, Fall 2026)
%% Shared Beamer preamble for every deck in slides/S01 ... slides/S10.
%%
%% Usage (a deck lives in slides/S0X/ and is compiled from that folder):
%%   %% ECON 282E -- Foundations of Macroeconomics (UC Riverside, Fall 2026)
%% Shared Beamer preamble for every deck in slides/S01 ... slides/S10.
%%
%% Usage (a deck lives in slides/S0X/ and is compiled from that folder):
%%   %% ECON 282E -- Foundations of Macroeconomics (UC Riverside, Fall 2026)
%% Shared Beamer preamble for every deck in slides/S01 ... slides/S10.
%%
%% Usage (a deck lives in slides/S0X/ and is compiled from that folder):
%%   \input{../preamble/course_preamble.tex}
%%   \decktitle{1}{A1}{Who I Am \& Why This Course}{September 24, 2026}
%%   \begin{document}
%%   \makedecktitle
%%   ...
%%
%% Adapted from Courses/AI-ECON-2026/source/shared_preamble.tex (Metropolis theme,
%% concept/caution/example/takeaway boxes, \sectiondivider). Changes for this course:
%%   * course title block (\decktitle / \makedecktitle) instead of \makebeamertitle
%%   * \zh{...} typesets Chinese (book titles) under pdflatex via CJKutf8
%%   * researchbox for the "Research in the wild" frames
%%   * section pages off; TikZ libraries and the course map loaded here
%% Build with pdflatex only (two passes); tools/build_slides.py does this for you.

\documentclass[10pt,english,aspectratio=169]{beamer}

\usepackage[T1]{fontenc}
\usepackage{textcomp}
\usepackage[utf8]{inputenc}
\usepackage{babel}
\usepackage{CJKutf8}

\usepackage{amsmath, amssymb, amsthm, graphicx}
\usepackage{booktabs, multirow, tabularx}
\usepackage{tikz}
\usetikzlibrary{positioning,arrows.meta,shapes,calc,fit,backgrounds}
\usepackage{pgfplots}
\pgfplotsset{compat=1.16}
\usepackage[most]{tcolorbox}
\tcbuselibrary{skins, breakable}
\usepackage{listings}
\usepackage{xcolor}

\lstset{
  basicstyle=\ttfamily\small,
  keywordstyle=\color{blue!80!black}\bfseries,
  commentstyle=\color{green!50!black}\itshape,
  stringstyle=\color{orange!80!black},
  backgroundcolor=\color{gray!8},
  frame=single,
  rulecolor=\color{gray!40},
  breaklines=true,
  showstringspaces=false,
  numbers=none,
}

\usetheme{metropolis}
\usefonttheme{professionalfonts}
\metroset{sectionpage=none}

\definecolor{LightGray}{RGB}{230,230,230}
\definecolor{RedTitle}{RGB}{0,0,200}
\definecolor{VeryLightGray}{RGB}{245,245,245}
\definecolor{DarkText}{RGB}{0,0,0}
\definecolor{UCRBlue}{RGB}{0,45,106}
\definecolor{UCRGold}{RGB}{241,171,0}

% Stage colours of the course arc (used by the course map and elsewhere)
\colorlet{StageTrunk}{blue!12}
\colorlet{StageBranch}{cyan!14}
\colorlet{StageMachine}{gray!18}
\colorlet{StageWall}{red!14}
\colorlet{StageTools}{orange!20}
\colorlet{StageData}{green!16}
\colorlet{StageAgents}{violet!14}

\hypersetup{bookmarks=false,breaklinks=true,pdfborder={0 0 0},colorlinks=true,
  linkcolor=DarkText,urlcolor=RedTitle}

\setbeamercolor{background canvas}{bg=white}
\setbeamercolor{title}{fg=RedTitle}
\setbeamercolor{frametitle}{bg=VeryLightGray,fg=DarkText}
\setbeamercolor{normal text}{fg=DarkText}
\setbeamercolor{alerted text}{fg=RedTitle}
\setbeamersize{text margin left=5mm, text margin right=5mm}
\addtobeamertemplate{frametitle}{}{\vspace{-0.5em}}

\setbeamertemplate{navigation symbols}{}
\setbeamertemplate{footline}{%
  \hspace{5mm}{\usebeamerfont{page number in head/foot}\color{black!45}%
  ECON 282E \,$\cdot$\, \insertshortdeck}\hfill%
  \usebeamercolor[fg]{page number in head/foot}%
  \usebeamerfont{page number in head/foot}%
  \insertframenumber\,/\,\inserttotalframenumber\kern1em\vskip2pt%
}

% ---------------------------------------------------------------------------
% Chinese text under pdflatex (book titles etc.)
\newcommand{\zh}[1]{\begin{CJK*}{UTF8}{gbsn}#1\end{CJK*}}

% ---------------------------------------------------------------------------
% Deck title block
%   \decktitle{session}{deck id}{title}{date}
\newcommand{\insertshortdeck}{}
\newcommand{\decksession}{}
\newcommand{\deckid}{}
\newcommand{\decktitletext}{}
\newcommand{\deckdate}{}
\newcommand{\decktitle}[4]{%
  \renewcommand{\decksession}{#1}%
  \renewcommand{\deckid}{#2}%
  \renewcommand{\decktitletext}{#3}%
  \renewcommand{\deckdate}{#4}%
  \renewcommand{\insertshortdeck}{Session #1 \,$\cdot$\, Deck #1.#2}%
  \title{#3}%
  \author{Zhigang Feng}%
  \date{#4}%
}
\newcommand{\makedecktitle}{%
  \begin{frame}[plain,noframenumbering]
    \begin{tikzpicture}[remember picture,overlay]
      \fill[UCRBlue] (current page.north west) rectangle ([yshift=-0.55cm]current page.north east);
      \fill[UCRGold] ([yshift=-0.55cm]current page.north west) rectangle ([yshift=-0.62cm]current page.north east);
      \node[anchor=west, text=white, font=\small\bfseries] at ([xshift=5mm,yshift=-0.275cm]current page.north west)
        {ECON 282E \,$\cdot$\, Foundations of Macroeconomics};
      \node[anchor=east, text=white, font=\small] at ([xshift=-5mm,yshift=-0.275cm]current page.north east)
        {UC Riverside \,$\cdot$\, Fall 2026};
    \end{tikzpicture}
    \vspace*{1.1cm}
    {\usebeamerfont{title}\color{black!55}\large Session \decksession\ \,$\cdot$\, Deck \decksession.\deckid\par}
    \vspace{0.25cm}
    {\usebeamerfont{title}\color{RedTitle}\LARGE\bfseries \decktitletext\par}
    \vspace{0.7cm}
    {\large Zhigang Feng\par}
    \vspace{0.15cm}
    {\color{black!65}\deckdate\par}
    \vfill
    {\footnotesize\itshape\color{black!55}Build it on paper $\;\to\;$ solve it on the machine $\;\to\;$ push it to the frontier with AI\par}
  \end{frame}}

% ---------------------------------------------------------------------------
% Section divider (plain frame, not counted)
\newcommand{\sectiondivider}[2]{%
\begin{frame}[plain,noframenumbering]
\begin{center}
\vfill
{\Large\textbf{#1}\par}
\vspace{0.35cm}
{\normalsize #2\par}
\vfill
\end{center}
\end{frame}
}

% ---------------------------------------------------------------------------
% Boxes
\newtcolorbox{conceptbox}[1][]{colback=blue!3, colframe=RedTitle!55, boxrule=0.35mm,
  arc=1mm, left=2mm, right=2mm, top=1mm, bottom=1mm, #1}
\newtcolorbox{cautionbox}[1][]{colback=red!3, colframe=red!50!black, boxrule=0.35mm,
  arc=1mm, left=2mm, right=2mm, top=1mm, bottom=1mm, #1}
\newtcolorbox{examplebox}[1][]{colback=green!3, colframe=green!45!black, boxrule=0.35mm,
  arc=1mm, left=2mm, right=2mm, top=1mm, bottom=1mm, #1}
\newtcolorbox{takeawaybox}[1][]{colback=VeryLightGray, colframe=RedTitle!55, boxrule=0.35mm,
  arc=1mm, left=2mm, right=2mm, top=1mm, bottom=1mm, #1}
\newtcolorbox{researchbox}[1][]{enhanced, colback=violet!4, colframe=violet!60!black,
  boxrule=0.35mm, arc=1mm, left=2mm, right=2mm, top=1mm, bottom=1mm,
  fonttitle=\bfseries\small, title={Research in the wild}, #1}

\newcommand{\compacttables}{\renewcommand{\arraystretch}{1.15}}
\newcommand{\src}[1]{{\tiny\color{black!55}#1}}

% ---------------------------------------------------------------------------
\input{../preamble/course_map.tex}

%%   \decktitle{1}{A1}{Who I Am \& Why This Course}{September 24, 2026}
%%   \begin{document}
%%   \makedecktitle
%%   ...
%%
%% Adapted from Courses/AI-ECON-2026/source/shared_preamble.tex (Metropolis theme,
%% concept/caution/example/takeaway boxes, \sectiondivider). Changes for this course:
%%   * course title block (\decktitle / \makedecktitle) instead of \makebeamertitle
%%   * \zh{...} typesets Chinese (book titles) under pdflatex via CJKutf8
%%   * researchbox for the "Research in the wild" frames
%%   * section pages off; TikZ libraries and the course map loaded here
%% Build with pdflatex only (two passes); tools/build_slides.py does this for you.

\documentclass[10pt,english,aspectratio=169]{beamer}

\usepackage[T1]{fontenc}
\usepackage{textcomp}
\usepackage[utf8]{inputenc}
\usepackage{babel}
\usepackage{CJKutf8}

\usepackage{amsmath, amssymb, amsthm, graphicx}
\usepackage{booktabs, multirow, tabularx}
\usepackage{tikz}
\usetikzlibrary{positioning,arrows.meta,shapes,calc,fit,backgrounds}
\usepackage{pgfplots}
\pgfplotsset{compat=1.16}
\usepackage[most]{tcolorbox}
\tcbuselibrary{skins, breakable}
\usepackage{listings}
\usepackage{xcolor}

\lstset{
  basicstyle=\ttfamily\small,
  keywordstyle=\color{blue!80!black}\bfseries,
  commentstyle=\color{green!50!black}\itshape,
  stringstyle=\color{orange!80!black},
  backgroundcolor=\color{gray!8},
  frame=single,
  rulecolor=\color{gray!40},
  breaklines=true,
  showstringspaces=false,
  numbers=none,
}

\usetheme{metropolis}
\usefonttheme{professionalfonts}
\metroset{sectionpage=none}

\definecolor{LightGray}{RGB}{230,230,230}
\definecolor{RedTitle}{RGB}{0,0,200}
\definecolor{VeryLightGray}{RGB}{245,245,245}
\definecolor{DarkText}{RGB}{0,0,0}
\definecolor{UCRBlue}{RGB}{0,45,106}
\definecolor{UCRGold}{RGB}{241,171,0}

% Stage colours of the course arc (used by the course map and elsewhere)
\colorlet{StageTrunk}{blue!12}
\colorlet{StageBranch}{cyan!14}
\colorlet{StageMachine}{gray!18}
\colorlet{StageWall}{red!14}
\colorlet{StageTools}{orange!20}
\colorlet{StageData}{green!16}
\colorlet{StageAgents}{violet!14}

\hypersetup{bookmarks=false,breaklinks=true,pdfborder={0 0 0},colorlinks=true,
  linkcolor=DarkText,urlcolor=RedTitle}

\setbeamercolor{background canvas}{bg=white}
\setbeamercolor{title}{fg=RedTitle}
\setbeamercolor{frametitle}{bg=VeryLightGray,fg=DarkText}
\setbeamercolor{normal text}{fg=DarkText}
\setbeamercolor{alerted text}{fg=RedTitle}
\setbeamersize{text margin left=5mm, text margin right=5mm}
\addtobeamertemplate{frametitle}{}{\vspace{-0.5em}}

\setbeamertemplate{navigation symbols}{}
\setbeamertemplate{footline}{%
  \hspace{5mm}{\usebeamerfont{page number in head/foot}\color{black!45}%
  ECON 282E \,$\cdot$\, \insertshortdeck}\hfill%
  \usebeamercolor[fg]{page number in head/foot}%
  \usebeamerfont{page number in head/foot}%
  \insertframenumber\,/\,\inserttotalframenumber\kern1em\vskip2pt%
}

% ---------------------------------------------------------------------------
% Chinese text under pdflatex (book titles etc.)
\newcommand{\zh}[1]{\begin{CJK*}{UTF8}{gbsn}#1\end{CJK*}}

% ---------------------------------------------------------------------------
% Deck title block
%   \decktitle{session}{deck id}{title}{date}
\newcommand{\insertshortdeck}{}
\newcommand{\decksession}{}
\newcommand{\deckid}{}
\newcommand{\decktitletext}{}
\newcommand{\deckdate}{}
\newcommand{\decktitle}[4]{%
  \renewcommand{\decksession}{#1}%
  \renewcommand{\deckid}{#2}%
  \renewcommand{\decktitletext}{#3}%
  \renewcommand{\deckdate}{#4}%
  \renewcommand{\insertshortdeck}{Session #1 \,$\cdot$\, Deck #1.#2}%
  \title{#3}%
  \author{Zhigang Feng}%
  \date{#4}%
}
\newcommand{\makedecktitle}{%
  \begin{frame}[plain,noframenumbering]
    \begin{tikzpicture}[remember picture,overlay]
      \fill[UCRBlue] (current page.north west) rectangle ([yshift=-0.55cm]current page.north east);
      \fill[UCRGold] ([yshift=-0.55cm]current page.north west) rectangle ([yshift=-0.62cm]current page.north east);
      \node[anchor=west, text=white, font=\small\bfseries] at ([xshift=5mm,yshift=-0.275cm]current page.north west)
        {ECON 282E \,$\cdot$\, Foundations of Macroeconomics};
      \node[anchor=east, text=white, font=\small] at ([xshift=-5mm,yshift=-0.275cm]current page.north east)
        {UC Riverside \,$\cdot$\, Fall 2026};
    \end{tikzpicture}
    \vspace*{1.1cm}
    {\usebeamerfont{title}\color{black!55}\large Session \decksession\ \,$\cdot$\, Deck \decksession.\deckid\par}
    \vspace{0.25cm}
    {\usebeamerfont{title}\color{RedTitle}\LARGE\bfseries \decktitletext\par}
    \vspace{0.7cm}
    {\large Zhigang Feng\par}
    \vspace{0.15cm}
    {\color{black!65}\deckdate\par}
    \vfill
    {\footnotesize\itshape\color{black!55}Build it on paper $\;\to\;$ solve it on the machine $\;\to\;$ push it to the frontier with AI\par}
  \end{frame}}

% ---------------------------------------------------------------------------
% Section divider (plain frame, not counted)
\newcommand{\sectiondivider}[2]{%
\begin{frame}[plain,noframenumbering]
\begin{center}
\vfill
{\Large\textbf{#1}\par}
\vspace{0.35cm}
{\normalsize #2\par}
\vfill
\end{center}
\end{frame}
}

% ---------------------------------------------------------------------------
% Boxes
\newtcolorbox{conceptbox}[1][]{colback=blue!3, colframe=RedTitle!55, boxrule=0.35mm,
  arc=1mm, left=2mm, right=2mm, top=1mm, bottom=1mm, #1}
\newtcolorbox{cautionbox}[1][]{colback=red!3, colframe=red!50!black, boxrule=0.35mm,
  arc=1mm, left=2mm, right=2mm, top=1mm, bottom=1mm, #1}
\newtcolorbox{examplebox}[1][]{colback=green!3, colframe=green!45!black, boxrule=0.35mm,
  arc=1mm, left=2mm, right=2mm, top=1mm, bottom=1mm, #1}
\newtcolorbox{takeawaybox}[1][]{colback=VeryLightGray, colframe=RedTitle!55, boxrule=0.35mm,
  arc=1mm, left=2mm, right=2mm, top=1mm, bottom=1mm, #1}
\newtcolorbox{researchbox}[1][]{enhanced, colback=violet!4, colframe=violet!60!black,
  boxrule=0.35mm, arc=1mm, left=2mm, right=2mm, top=1mm, bottom=1mm,
  fonttitle=\bfseries\small, title={Research in the wild}, #1}

\newcommand{\compacttables}{\renewcommand{\arraystretch}{1.15}}
\newcommand{\src}[1]{{\tiny\color{black!55}#1}}

% ---------------------------------------------------------------------------
%% Course map: the recurring "you are here" slide for ECON 282E.
%% Loaded by course_preamble.tex. Pattern follows AI-ECON-2026 lec01_map.tex, but the map
%% shows the whole ten-session course rather than one lecture.
%%
%%   \CourseMap{s}{label}   s = 1..10 highlights session s and prints "you are here: label"
%%                          (e.g. \CourseMap{1}{1.A2}); earlier sessions get a check mark.
%%   \CourseMap{0}{}        full overview, nothing highlighted.
%%
%% Note: TikZ option lists are comma-split before TeX conditionals run, so every \ifnum
%% below sits outside [...] option lists.

\newcommand{\CourseMap}[2]{%
\begin{center}
\begin{tikzpicture}[
    sess/.style={rectangle, rounded corners=2pt, draw=black!45, minimum width=1.34cm,
        minimum height=1.12cm, text width=1.26cm, align=center, inner sep=1pt},
    stage/.style={font=\tiny\bfseries, text=black!70},
]
  \foreach \i/\d/\t/\c in {%
      1/Sep 24/Intro \\ Growth/StageTrunk,
      2/Oct 1/HA $\cdot$ OLG \\ Policy/StageBranch,
      3/Oct 8/Num.\ DP \\ Python/StageMachine,
      4/Oct 15/Numerics \\ I/StageMachine,
      5/Oct 22/Numerics \\ II/StageMachine,
      6/Oct 29/The wall \\ \textbf{Midterm}/StageWall,
      7/Nov 5/ML \\ Deep nets/StageTools,
      8/Nov 12/RL \\ HA $+$ DL/StageTools,
      9/Nov 19/Text \\ LLMs/StageData,
      10/Dec 3/Agents \\ \textbf{Projects}/StageAgents} {
    \node[sess, fill=\c] (n\i) at ({1.46*(\i-1)}, 0)
      {{\scriptsize\bfseries S\i}\\[-1pt]{\tiny\color{black!60}\d}\\[1pt]{\tiny \t}};
    \ifnum\i<#1
      \node[font=\scriptsize, text=green!45!black] at ([xshift=-2.5mm,yshift=-2.2mm]n\i.north east) {$\checkmark$};
    \fi
    \ifnum\i=#1
      \draw[RedTitle, very thick, rounded corners=2pt]
        ([xshift=-0.6mm,yshift=-0.6mm]n\i.south west) rectangle ([xshift=0.6mm,yshift=0.6mm]n\i.north east);
      % keep the label inside the picture at both ends of the row
      \ifnum\i<3
        \node[font=\scriptsize\bfseries, text=RedTitle, anchor=south west] at ([yshift=1.2mm]n\i.north west)
          {$\blacktriangledown$\,you are here: #2};
      \else\ifnum\i>8
        \node[font=\scriptsize\bfseries, text=RedTitle, anchor=south east] at ([yshift=1.2mm]n\i.north east)
          {you are here: #2\,$\blacktriangledown$};
      \else
        \node[font=\scriptsize\bfseries, text=RedTitle, anchor=south] at ([yshift=1.2mm]n\i.north)
          {you are here: #2\,$\blacktriangledown$};
      \fi\fi
    \fi
  }
  % stage band under the sessions
  \foreach \a/\b/\lab/\c in {1/1/trunk/StageTrunk, 2/2/branches/StageBranch,
      3/5/the machine $\cdot$ classical methods/StageMachine, 6/6/the wall/StageWall,
      7/8/new tools/StageTools, 9/9/new data/StageData, 10/10/colleagues/StageAgents} {
    \fill[\c] ([yshift=-1.5mm]n\a.south west) rectangle ([yshift=-4.6mm]n\b.south east);
    \path ([yshift=-1.5mm]n\a.south west) -- ([yshift=-4.6mm]n\b.south east)
      node[midway, stage] {\lab};
  }
  \node[font=\footnotesize\itshape, text=black!70] at ([yshift=-9.5mm]$(n1.south)!0.5!(n10.south)$)
    {Build it on paper $\;\to\;$ solve it on the machine $\;\to\;$ push it to the frontier with AI};
\end{tikzpicture}
\end{center}
}


%%   \decktitle{1}{A1}{Who I Am \& Why This Course}{September 24, 2026}
%%   \begin{document}
%%   \makedecktitle
%%   ...
%%
%% Adapted from Courses/AI-ECON-2026/source/shared_preamble.tex (Metropolis theme,
%% concept/caution/example/takeaway boxes, \sectiondivider). Changes for this course:
%%   * course title block (\decktitle / \makedecktitle) instead of \makebeamertitle
%%   * \zh{...} typesets Chinese (book titles) under pdflatex via CJKutf8
%%   * researchbox for the "Research in the wild" frames
%%   * section pages off; TikZ libraries and the course map loaded here
%% Build with pdflatex only (two passes); tools/build_slides.py does this for you.

\documentclass[10pt,english,aspectratio=169]{beamer}

\usepackage[T1]{fontenc}
\usepackage{textcomp}
\usepackage[utf8]{inputenc}
\usepackage{babel}
\usepackage{CJKutf8}

\usepackage{amsmath, amssymb, amsthm, graphicx}
\usepackage{booktabs, multirow, tabularx}
\usepackage{tikz}
\usetikzlibrary{positioning,arrows.meta,shapes,calc,fit,backgrounds}
\usepackage{pgfplots}
\pgfplotsset{compat=1.16}
\usepackage[most]{tcolorbox}
\tcbuselibrary{skins, breakable}
\usepackage{listings}
\usepackage{xcolor}

\lstset{
  basicstyle=\ttfamily\small,
  keywordstyle=\color{blue!80!black}\bfseries,
  commentstyle=\color{green!50!black}\itshape,
  stringstyle=\color{orange!80!black},
  backgroundcolor=\color{gray!8},
  frame=single,
  rulecolor=\color{gray!40},
  breaklines=true,
  showstringspaces=false,
  numbers=none,
}

\usetheme{metropolis}
\usefonttheme{professionalfonts}
\metroset{sectionpage=none}

\definecolor{LightGray}{RGB}{230,230,230}
\definecolor{RedTitle}{RGB}{0,0,200}
\definecolor{VeryLightGray}{RGB}{245,245,245}
\definecolor{DarkText}{RGB}{0,0,0}
\definecolor{UCRBlue}{RGB}{0,45,106}
\definecolor{UCRGold}{RGB}{241,171,0}

% Stage colours of the course arc (used by the course map and elsewhere)
\colorlet{StageTrunk}{blue!12}
\colorlet{StageBranch}{cyan!14}
\colorlet{StageMachine}{gray!18}
\colorlet{StageWall}{red!14}
\colorlet{StageTools}{orange!20}
\colorlet{StageData}{green!16}
\colorlet{StageAgents}{violet!14}

\hypersetup{bookmarks=false,breaklinks=true,pdfborder={0 0 0},colorlinks=true,
  linkcolor=DarkText,urlcolor=RedTitle}

\setbeamercolor{background canvas}{bg=white}
\setbeamercolor{title}{fg=RedTitle}
\setbeamercolor{frametitle}{bg=VeryLightGray,fg=DarkText}
\setbeamercolor{normal text}{fg=DarkText}
\setbeamercolor{alerted text}{fg=RedTitle}
\setbeamersize{text margin left=5mm, text margin right=5mm}
\addtobeamertemplate{frametitle}{}{\vspace{-0.5em}}

\setbeamertemplate{navigation symbols}{}
\setbeamertemplate{footline}{%
  \hspace{5mm}{\usebeamerfont{page number in head/foot}\color{black!45}%
  ECON 282E \,$\cdot$\, \insertshortdeck}\hfill%
  \usebeamercolor[fg]{page number in head/foot}%
  \usebeamerfont{page number in head/foot}%
  \insertframenumber\,/\,\inserttotalframenumber\kern1em\vskip2pt%
}

% ---------------------------------------------------------------------------
% Chinese text under pdflatex (book titles etc.)
\newcommand{\zh}[1]{\begin{CJK*}{UTF8}{gbsn}#1\end{CJK*}}

% ---------------------------------------------------------------------------
% Deck title block
%   \decktitle{session}{deck id}{title}{date}
\newcommand{\insertshortdeck}{}
\newcommand{\decksession}{}
\newcommand{\deckid}{}
\newcommand{\decktitletext}{}
\newcommand{\deckdate}{}
\newcommand{\decktitle}[4]{%
  \renewcommand{\decksession}{#1}%
  \renewcommand{\deckid}{#2}%
  \renewcommand{\decktitletext}{#3}%
  \renewcommand{\deckdate}{#4}%
  \renewcommand{\insertshortdeck}{Session #1 \,$\cdot$\, Deck #1.#2}%
  \title{#3}%
  \author{Zhigang Feng}%
  \date{#4}%
}
\newcommand{\makedecktitle}{%
  \begin{frame}[plain,noframenumbering]
    \begin{tikzpicture}[remember picture,overlay]
      \fill[UCRBlue] (current page.north west) rectangle ([yshift=-0.55cm]current page.north east);
      \fill[UCRGold] ([yshift=-0.55cm]current page.north west) rectangle ([yshift=-0.62cm]current page.north east);
      \node[anchor=west, text=white, font=\small\bfseries] at ([xshift=5mm,yshift=-0.275cm]current page.north west)
        {ECON 282E \,$\cdot$\, Foundations of Macroeconomics};
      \node[anchor=east, text=white, font=\small] at ([xshift=-5mm,yshift=-0.275cm]current page.north east)
        {UC Riverside \,$\cdot$\, Fall 2026};
    \end{tikzpicture}
    \vspace*{1.1cm}
    {\usebeamerfont{title}\color{black!55}\large Session \decksession\ \,$\cdot$\, Deck \decksession.\deckid\par}
    \vspace{0.25cm}
    {\usebeamerfont{title}\color{RedTitle}\LARGE\bfseries \decktitletext\par}
    \vspace{0.7cm}
    {\large Zhigang Feng\par}
    \vspace{0.15cm}
    {\color{black!65}\deckdate\par}
    \vfill
    {\footnotesize\itshape\color{black!55}Build it on paper $\;\to\;$ solve it on the machine $\;\to\;$ push it to the frontier with AI\par}
  \end{frame}}

% ---------------------------------------------------------------------------
% Section divider (plain frame, not counted)
\newcommand{\sectiondivider}[2]{%
\begin{frame}[plain,noframenumbering]
\begin{center}
\vfill
{\Large\textbf{#1}\par}
\vspace{0.35cm}
{\normalsize #2\par}
\vfill
\end{center}
\end{frame}
}

% ---------------------------------------------------------------------------
% Boxes
\newtcolorbox{conceptbox}[1][]{colback=blue!3, colframe=RedTitle!55, boxrule=0.35mm,
  arc=1mm, left=2mm, right=2mm, top=1mm, bottom=1mm, #1}
\newtcolorbox{cautionbox}[1][]{colback=red!3, colframe=red!50!black, boxrule=0.35mm,
  arc=1mm, left=2mm, right=2mm, top=1mm, bottom=1mm, #1}
\newtcolorbox{examplebox}[1][]{colback=green!3, colframe=green!45!black, boxrule=0.35mm,
  arc=1mm, left=2mm, right=2mm, top=1mm, bottom=1mm, #1}
\newtcolorbox{takeawaybox}[1][]{colback=VeryLightGray, colframe=RedTitle!55, boxrule=0.35mm,
  arc=1mm, left=2mm, right=2mm, top=1mm, bottom=1mm, #1}
\newtcolorbox{researchbox}[1][]{enhanced, colback=violet!4, colframe=violet!60!black,
  boxrule=0.35mm, arc=1mm, left=2mm, right=2mm, top=1mm, bottom=1mm,
  fonttitle=\bfseries\small, title={Research in the wild}, #1}

\newcommand{\compacttables}{\renewcommand{\arraystretch}{1.15}}
\newcommand{\src}[1]{{\tiny\color{black!55}#1}}

% ---------------------------------------------------------------------------
%% Course map: the recurring "you are here" slide for ECON 282E.
%% Loaded by course_preamble.tex. Pattern follows AI-ECON-2026 lec01_map.tex, but the map
%% shows the whole ten-session course rather than one lecture.
%%
%%   \CourseMap{s}{label}   s = 1..10 highlights session s and prints "you are here: label"
%%                          (e.g. \CourseMap{1}{1.A2}); earlier sessions get a check mark.
%%   \CourseMap{0}{}        full overview, nothing highlighted.
%%
%% Note: TikZ option lists are comma-split before TeX conditionals run, so every \ifnum
%% below sits outside [...] option lists.

\newcommand{\CourseMap}[2]{%
\begin{center}
\begin{tikzpicture}[
    sess/.style={rectangle, rounded corners=2pt, draw=black!45, minimum width=1.34cm,
        minimum height=1.12cm, text width=1.26cm, align=center, inner sep=1pt},
    stage/.style={font=\tiny\bfseries, text=black!70},
]
  \foreach \i/\d/\t/\c in {%
      1/Sep 24/Intro \\ Growth/StageTrunk,
      2/Oct 1/HA $\cdot$ OLG \\ Policy/StageBranch,
      3/Oct 8/Num.\ DP \\ Python/StageMachine,
      4/Oct 15/Numerics \\ I/StageMachine,
      5/Oct 22/Numerics \\ II/StageMachine,
      6/Oct 29/The wall \\ \textbf{Midterm}/StageWall,
      7/Nov 5/ML \\ Deep nets/StageTools,
      8/Nov 12/RL \\ HA $+$ DL/StageTools,
      9/Nov 19/Text \\ LLMs/StageData,
      10/Dec 3/Agents \\ \textbf{Projects}/StageAgents} {
    \node[sess, fill=\c] (n\i) at ({1.46*(\i-1)}, 0)
      {{\scriptsize\bfseries S\i}\\[-1pt]{\tiny\color{black!60}\d}\\[1pt]{\tiny \t}};
    \ifnum\i<#1
      \node[font=\scriptsize, text=green!45!black] at ([xshift=-2.5mm,yshift=-2.2mm]n\i.north east) {$\checkmark$};
    \fi
    \ifnum\i=#1
      \draw[RedTitle, very thick, rounded corners=2pt]
        ([xshift=-0.6mm,yshift=-0.6mm]n\i.south west) rectangle ([xshift=0.6mm,yshift=0.6mm]n\i.north east);
      % keep the label inside the picture at both ends of the row
      \ifnum\i<3
        \node[font=\scriptsize\bfseries, text=RedTitle, anchor=south west] at ([yshift=1.2mm]n\i.north west)
          {$\blacktriangledown$\,you are here: #2};
      \else\ifnum\i>8
        \node[font=\scriptsize\bfseries, text=RedTitle, anchor=south east] at ([yshift=1.2mm]n\i.north east)
          {you are here: #2\,$\blacktriangledown$};
      \else
        \node[font=\scriptsize\bfseries, text=RedTitle, anchor=south] at ([yshift=1.2mm]n\i.north)
          {you are here: #2\,$\blacktriangledown$};
      \fi\fi
    \fi
  }
  % stage band under the sessions
  \foreach \a/\b/\lab/\c in {1/1/trunk/StageTrunk, 2/2/branches/StageBranch,
      3/5/the machine $\cdot$ classical methods/StageMachine, 6/6/the wall/StageWall,
      7/8/new tools/StageTools, 9/9/new data/StageData, 10/10/colleagues/StageAgents} {
    \fill[\c] ([yshift=-1.5mm]n\a.south west) rectangle ([yshift=-4.6mm]n\b.south east);
    \path ([yshift=-1.5mm]n\a.south west) -- ([yshift=-4.6mm]n\b.south east)
      node[midway, stage] {\lab};
  }
  \node[font=\footnotesize\itshape, text=black!70] at ([yshift=-9.5mm]$(n1.south)!0.5!(n10.south)$)
    {Build it on paper $\;\to\;$ solve it on the machine $\;\to\;$ push it to the frontier with AI};
\end{tikzpicture}
\end{center}
}


\graphicspath{{./figures/}}

\lstset{language=Python, basicstyle=\ttfamily\scriptsize, frame=none,
        backgroundcolor=\color{gray!8}, aboveskip=2pt, belowskip=2pt}

\decktitle{4}{A3}{Optimization in Many Dimensions}{Thursday, October 15, 2026}

\begin{document}

\makedecktitle

%=============================================================================
\begin{frame}{Where We Are}
\CourseMap{4}{4.A3}
\vspace{-1mm}
\begin{center}
\small \textbf{Block A:} \textbf{4.A1} root-finding $\;\cdot\;$ \textbf{4.A2} optimization on an interval $\;\cdot\;$ \textbf{4.A3} optimization in many dimensions $\;\cdot\;$ \textbf{4.A4} differentiation.
\end{center}
\end{frame}

%=============================================================================
\begin{frame}{What Breaks When the Choice Is a Vector}
\begin{columns}[T]
\begin{column}{0.51\textwidth}
\small
Every method in 4.A2 worked by \textbf{shrinking a bracket}. An interval has two ends and a clear inside. A region of $\mathbb{R}^{N}$ has neither.
\medskip

What replaces it is a \textbf{local quadratic model}:
\[
  f(\theta+\Delta)\approx f(\theta)+g^{\!\top}\!\Delta+\tfrac12\Delta^{\!\top} H\Delta,
\]
with $g=\nabla f$ and $H=\nabla^{2}f$. Setting its gradient to zero gives the Newton step
\[
  \Delta=-H^{-1}g .
\]
\end{column}
\begin{column}{0.46\textwidth}
\begin{center}\footnotesize
\begin{tabularx}{\linewidth}{@{}c r r@{}}
\toprule
$N$ & \textbf{Hessian entries} & \textbf{flops to solve} \\
\midrule
10 & 100 & 333 \\
100 & 10{,}000 & 333{,}333 \\
1{,}000 & \textbf{1 million} & \textbf{333 million} \\
\bottomrule
\end{tabularx}
\end{center}
\vspace{1mm}
\begin{cautionbox}\footnotesize
And that is \emph{per step}. Storing $H$ is $O(N^{2})$, factorising it $O(N^{3})$, and building it by finite differences costs $O(N^{2})$ function evaluations on top.
\end{cautionbox}
\end{column}
\end{columns}
\vspace{1mm}
\src{Cost table adapted from QM \emph{Numerical Methods}, ``Why Multivariate is Hard''.}
\end{frame}

%=============================================================================
\begin{frame}{Quasi-Newton: Learn the Curvature From the Steps You Took}
\begin{columns}[T]
\begin{column}{0.53\textwidth}
\small
You already know something about $H$ after one step. If you moved by $s=\theta_1-\theta_0$ and the gradient changed by $y=g_1-g_0$, then for a quadratic
\[
  H\,s=y .
\]
That is the \textbf{secant condition}. Quasi-Newton methods carry an approximation $B$ and update it so that $B_{1}s_{0}=y_{0}$ holds exactly, while changing $B$ as little as possible.
\medskip

BFGS is the rank-two update that does this and keeps $B$ positive definite:
\[
  B_{1}=B_{0}+\frac{y_0y_0^{\!\top}}{y_0^{\!\top}s_0}-\frac{B_0s_0s_0^{\!\top}B_0}{s_0^{\!\top}B_0s_0}.
\]
\end{column}
\begin{column}{0.44\textwidth}
\begin{conceptbox}[title=\textbf{What it buys}]\small
No second derivatives, ever. The per-step cost falls from $O(N^{3})$ to $O(N^{2})$, and L-BFGS --- which stores only the last few $(s,y)$ pairs --- brings it to $O(N)$.
\end{conceptbox}
\medskip
\begin{examplebox}\footnotesize
This is the same trade as Broyden in 4.A1: pay a few more iterations, stop building derivative matrices.
\end{examplebox}
\end{column}
\end{columns}
\end{frame}

%=============================================================================
\begin{frame}{BFGS, Worked Through by Hand}
\begin{columns}[T]
\begin{column}{0.50\textwidth}
\small
Minimize $f(x,y)=\tfrac12x^{2}+y^{2}$ from $\theta_0=(2,2)$, so $g_0=(2,4)$ and $H=\mathrm{diag}(1,2)$.
\medskip
\begin{enumerate}\footnotesize
  \item $B_0=I$, so $d_0=-g_0=(-2,-4)$.
  \item Step $\alpha_0=0.5$: \;$\theta_1=(1,0)$, \;$g_1=(1,0)$.
  \item $s_0=(-1,-2)$, \;$y_0=(-1,-4)$.
  \item $y_0^{\!\top}s_0=9$, \;$s_0^{\!\top}B_0s_0=5$.
  \item $B_1=\begin{pmatrix}41/45 & 2/45\\[1pt] 2/45 & 89/45\end{pmatrix}
        =\begin{pmatrix}0.911 & 0.044\\[1pt] 0.044 & 1.978\end{pmatrix}$.
\end{enumerate}
\end{column}
\begin{column}{0.47\textwidth}
\begin{takeawaybox}[title=\textbf{Check the secant condition}]\small
$B_1s_0-y_0$ has largest entry $4\times10^{-16}$ --- it holds \textbf{exactly}, not approximately.
\end{takeawaybox}
\medskip
\begin{cautionbox}\footnotesize
Sources that print $B_1$ rounded to two decimals and then report ``$B_1s_0\approx y_0$'' are describing their own rounding, not the method. The condition is imposed by construction.
\end{cautionbox}
\medskip
\begin{examplebox}\footnotesize
Sanity check: the true Hessian gives $Hs_0=(-1,-4)=y_0$. One step has already taught $B$ the curvature along $s_0$.
\end{examplebox}
\end{column}
\end{columns}
\vspace{1mm}
\src{Adapted from QM \emph{Numerical Methods}, BFGS example. The step $\alpha_0=0.5$ is a convenient round number; the exact line-search minimizer is $5/9$.}
\end{frame}

%=============================================================================
\begin{frame}{The Step Size: What a Line Search Is For}
\begin{columns}[T]
\begin{column}{0.52\textwidth}
\small
The update gives a \emph{direction} $d$; how far to move along it is a separate, one-dimensional problem:
\[
  \varphi(\alpha)=f(\theta+\alpha d), \qquad \alpha>0 .
\]
Solving it exactly is wasteful --- the direction will change next step anyway. What is needed is a step that makes \textbf{enough} progress:
\begin{itemize}\footnotesize
  \item \textbf{Armijo}: $f$ falls by a fixed fraction of what the slope promised.
  \item \textbf{Curvature}: the slope has flattened enough that we are not stopping too early.
\end{itemize}
Together these are the \textbf{Wolfe conditions}; backtracking from $\alpha=1$ is the usual implementation.
\end{column}
\begin{column}{0.45\textwidth}
\begin{conceptbox}[title=\textbf{This is 4.A2, nested}]\small
The inner problem is exactly the scalar maximization of the previous deck --- which is why golden section and Brent turn up inside multivariate optimizers.
\end{conceptbox}
\medskip
\begin{cautionbox}\footnotesize
The curvature condition is not optional for BFGS: it is what guarantees $y^{\!\top}s>0$, and hence that $B$ stays positive definite. Skip it and the update can produce a direction that is not even uphill.
\end{cautionbox}
\end{column}
\end{columns}
\vspace{1mm}
\src{Adapted from QM \emph{Numerical Methods}, ``Interlude: What is a Line Search?''.}
\end{frame}

%=============================================================================
\begin{frame}{Conjugate Gradient: $N$ Steps, and No Matrix at All}
\begin{columns}[T]
\begin{column}{0.53\textwidth}
\small
Steepest descent zig-zags: each step undoes part of the last, because consecutive directions are orthogonal rather than \emph{complementary}.
\medskip

Conjugate gradient fixes the direction instead of the step. Choose $d_n$ \textbf{$H$-conjugate} to everything already used,
\[
  d_i^{\!\top}Hd_j=0 \quad (i\neq j),
\]
by the Fletcher--Reeves recursion
\[
  d_{n+1}=-g_{n+1}+\frac{g_{n+1}^{\!\top}g_{n+1}}{g_n^{\!\top}g_n}\,d_n .
\]
Each step then removes one direction's error \textbf{permanently}.
\end{column}
\begin{column}{0.44\textwidth}
\begin{takeawaybox}[title=\textbf{Exact in $N$ steps}]\small
On our quadratic, from $(2,2)$:
\[
  (2,2)\to(0.889,-0.222)\to(0,0),
\]
converged to machine zero in \textbf{two} steps in two dimensions --- as the theory promises.
\end{takeawaybox}
\medskip
\begin{examplebox}\footnotesize
And $H$ never appears. CG needs only the ability to \emph{multiply} by $H$ --- or, in the non-linear version, nothing but gradients.
\end{examplebox}
\end{column}
\end{columns}
\vspace{1mm}
\src{Adapted from QM Lec.~4's conjugate-gradient block. Here $d_1=[-80/81,\,20/81]$; the next two steps use that value.}
\end{frame}

%=============================================================================
\begin{frame}{The Memory Wall}
\begin{columns}[T]
\begin{column}{0.52\textwidth}
\small
The reason CG exists is not elegance. It is that $B$ does not fit.
\medskip
\begin{center}\footnotesize
\begin{tabularx}{\linewidth}{@{}c r r r@{}}
\toprule
$N$ & \textbf{Newton / BFGS} & \textbf{L-BFGS} & \textbf{CG} \\
\midrule
$10^{2}$ & $10^{4}$ & $2{,}000$ & $300$ \\
$10^{4}$ & $10^{8}$ & $2\times10^{5}$ & $3\times10^{4}$ \\
$10^{6}$ & $\mathbf{10^{12}}$ & $2\times10^{7}$ & $3\times10^{6}$ \\
\bottomrule
\end{tabularx}
\end{center}
\medskip
At $N=10^{6}$ the dense approximation is $10^{12}$ doubles --- \textbf{8 terabytes}. L-BFGS needs 160\,MB and CG needs 24.
\end{column}
\begin{column}{0.45\textwidth}
\begin{conceptbox}[title=\textbf{L-BFGS}]\small
Keep the last $m\approx10$ pairs $(s,y)$ and rebuild the action of $B^{-1}$ on demand. $O(mN)$ memory, and almost all of BFGS's behaviour.
\end{conceptbox}
\medskip
\begin{takeawaybox}\footnotesize
$N=10^{6}$ is not hypothetical. It is the parameter count of a small neural network --- which is why \textbf{no one trains one with Newton's method}, and why Session 7's optimizer is a first-order one.
\end{takeawaybox}
\end{column}
\end{columns}
\end{frame}

%=============================================================================
\begin{frame}{When the Objective Is a Sum Over Data}
\begin{columns}[T]
\begin{column}{0.53\textwidth}
\small
A different kind of large problem. Suppose
\[
  L(\theta)=\frac1M\sum_{i=1}^{M}\ell(\theta;z_i),
\]
with $M$ in the millions --- an empirical risk, a simulated method of moments objective, a likelihood over a long panel.
\medskip

Here the trouble is not the dimension of $\theta$ but the \textbf{cost of one gradient}: every step reads all $M$ observations, so a single BFGS iteration is a full pass over the data.
\medskip

\textbf{Stochastic gradient descent} replaces the full gradient with the gradient of a random subset:
\[
  \theta_{n+1}=\theta_n-\alpha_n\nabla_\theta\ell(\theta_n;z_{i_n}).
\]
\end{column}
\begin{column}{0.44\textwidth}
\begin{conceptbox}[title=\textbf{Why it is not madness}]\small
The sampled gradient is \textbf{unbiased}: its expectation over the draw is the full gradient. You are taking a noisy step in a direction that is right on average.
\end{conceptbox}
\medskip
\begin{examplebox}\footnotesize
The trade is stark --- $M$ times cheaper per step, and a step that is wrong. Whether that is a good bargain depends entirely on how much the observations agree with each other.
\end{examplebox}
\end{column}
\end{columns}
\vspace{1mm}
\src{Adapted from QM Lec.~4's SGD block.}
\end{frame}

%=============================================================================
\begin{frame}{The Noisy Path, and the Dimension Penalty}
\begin{columns}[T]
\begin{column}{0.51\textwidth}
\small
Split our $f=\tfrac12x^{2}+y^{2}$ into $\ell_1=\tfrac12x^{2}$ and $\ell_2=y^{2}$, and draw one at a time with $\alpha=0.5$:
\medskip
\begin{center}\scriptsize
\begin{tabularx}{\linewidth}{@{}>{\raggedright\arraybackslash}X c >{\raggedright\arraybackslash}X@{}}
\toprule
$\theta_n$ & \textbf{drawn} & $\theta_{n+1}$ \\
\midrule
$(2,\,2)$ & $\ell_2$ & $(2,\,0)$ \\
$(2,\,0)$ & $\ell_1$ & $(1,\,0)$ \\
$(1,\,0)$ & $\ell_1$ & $(0.5,\,0)$ \\
$(0.5,\,0)$ & $\ell_2$ & $(0.5,\,0)$ \\
$(0.5,\,0)$ & $\ell_1$ & $(0.25,\,0)$ \\
\bottomrule
\end{tabularx}
\end{center}
\medskip
\footnotesize
Row four makes no progress at all --- the draw carried no information about $x$. That is the noise, and it is the price of the cheap step.
\end{column}
\begin{column}{0.46\textwidth}
\begin{cautionbox}[title=\textbf{Why $\alpha$ must shrink with $N$}]\footnotesize
Noise of unit size in each coordinate has norm $\sqrt N$:
\medskip
\begin{center}\scriptsize
\begin{tabular}{@{}c r@{}}
\toprule
$N$ & $\lVert\text{noise}\rVert$ \\
\midrule
1 & 1 \\
$10^{4}$ & 100 \\
$10^{6}$ & \textbf{1{,}000} \\
\bottomrule
\end{tabular}
\end{center}
\medskip
so a step safe in one dimension is a thousand times too large in a million. Convergence needs $\sum\alpha_n=\infty$ and $\sum\alpha_n^{2}<\infty$ --- Robbins--Monro.
\end{cautionbox}
\end{column}
\end{columns}
\vspace{1mm}
\src{Trajectory from QM Lec.~4's SGD example, recomputed. \textbf{Note:} the source defines $f$ as the \emph{sum} and then differentiates it with $1/M$ averaging, giving $\nabla f(2,2)=[1,2]$ --- which contradicts the $[2,4]$ its BFGS and CG examples use for the same function, and contradicts this table. The consistent definition is used here.}
\end{frame}

%=============================================================================
\begin{frame}{When You Have No Derivatives at All}
\begin{columns}[T]
\begin{column}{0.51\textwidth}
\small
\textbf{Nelder--Mead} keeps a simplex of $N+1$ points and moves it by reflecting the worst vertex through the centroid, then expanding, contracting or shrinking. No derivatives, no model of $f$, and no convergence theory worth the name --- but it is robust to noise and kinks.
\medskip
\begin{center}\footnotesize
\begin{tabularx}{\linewidth}{@{}>{\raggedright\arraybackslash}X r r@{}}
\toprule
Rosenbrock from $(-1.2,1)$ & \textbf{evals} & \textbf{$|\theta-\theta^*|$} \\
\midrule
BFGS & \textbf{117} & $9\times10^{-6}$ \\
Nelder--Mead & 249 & $5\times10^{-11}$ \\
Powell & 607 & $2\times10^{-13}$ \\
\bottomrule
\end{tabularx}
\end{center}
\end{column}
\begin{column}{0.46\textwidth}
\begin{takeawaybox}[title=\textbf{Read that carefully}]\small
BFGS is fastest by a factor of two, and \emph{least} accurate by five orders of magnitude --- it stops when its gradient test is satisfied, which on a curved valley happens early.
\end{takeawaybox}
\medskip
\begin{examplebox}\footnotesize
``Fastest'' and ``most accurate'' are different questions. Decide which one you are asking \textbf{before} you pick the method, and report the tolerance you used.
\end{examplebox}
\end{column}
\end{columns}
\end{frame}

%=============================================================================
\begin{frame}{Simulated Annealing: Accept a Worse Point, On Purpose}
\begin{columns}[T]
\begin{column}{0.52\textwidth}
\small
Nelder--Mead needs no derivatives but still walks downhill, so it still stops at the first optimum it meets. When the objective is \textbf{multimodal}, that is the wrong answer, and no amount of accuracy fixes it.
\medskip

Simulated annealing proposes a random move and accepts it with probability
\[
  \min\Big\{1,\ \exp\big(\Delta f/T_n\big)\Big\},
\]
so an \emph{uphill} move is always taken and a downhill one is sometimes taken --- often when $T$ is high, rarely as $T_n\to0$.
\medskip
\footnotesize
Cooling slowly enough ($T_n\sim c/\ln n$) gives convergence to the global optimum in probability.
\end{column}
\begin{column}{0.45\textwidth}
\begin{center}\footnotesize
\begin{tabularx}{\linewidth}{@{}>{\raggedright\arraybackslash}X r r@{}}
\toprule
from $x_0=8$ & $x$ & $f(x)$ \\
\midrule
local method & 8.35 & $-0.50$ \\
annealing & \textbf{0.00} & $\mathbf{3.00}$ \\
\midrule
truth & 0.00 & 3.00 \\
\bottomrule
\end{tabularx}
\end{center}
\vspace{1mm}
\begin{takeawaybox}\footnotesize
On $-x^{2}/20+3\cos 3x$ the local method is trapped in the valley it started in. Annealing walks out of it and finds the global maximum --- in 20{,}000 evaluations, which is the price.
\end{takeawaybox}
\end{column}
\end{columns}
\vspace{1mm}
\src{Adapted from QM Lec.~4, ``Simulated Annealing: MCMC for Optimization''.}
\end{frame}

%=============================================================================
\begin{frame}{One Accept--Reject Rule, Two Very Different Uses}
\begin{columns}[T]
\begin{column}{0.53\textwidth}
\small
The rule on the last slide is not really an optimization idea. It is the \textbf{Metropolis--Hastings} rule, and the only thing that changes between its two uses is what you do with the temperature.
\medskip
\begin{center}\footnotesize
\begin{tabularx}{\linewidth}{@{}>{\raggedright\arraybackslash}p{1.6cm} >{\raggedright\arraybackslash}X >{\raggedright\arraybackslash}X@{}}
\toprule
 & \textbf{annealing} & \textbf{Metropolis--Hastings} \\
\midrule
accept w.p. & $\min\{1,e^{\Delta f/T_n}\}$ & $\min\{1,\tfrac{\pi(\theta^{*})}{\pi(\theta_n)}\}$ \\
temperature & driven to $0$ & held at $1$ \\
you want & the \textbf{argmax} & the \textbf{distribution} \\
output & one point & a whole chain \\
\bottomrule
\end{tabularx}
\end{center}
\end{column}
\begin{column}{0.44\textwidth}
\begin{conceptbox}[title=\textbf{Why an economist cares}]\small
Hold the temperature at one and the chain's stationary distribution \emph{is} the posterior. That is how a structural model is estimated when the likelihood can be evaluated but not maximised reliably.
\end{conceptbox}
\medskip
\begin{takeawaybox}\footnotesize
We use the optimization half today. The sampling half returns in \textbf{Session 9}, where calibration and estimation are the bridge from models to data --- and where the object being sampled is a posterior over parameters rather than a maximum.
\end{takeawaybox}
\end{column}
\end{columns}
\vspace{1mm}
\src{Comparison adapted from QM Lec.~4, ``RWMH vs.\ Simulated Annealing'' and ``When to Use?''.}
\end{frame}

%=============================================================================
\begin{frame}{When Constraints Bind}
\begin{columns}[T]
\begin{column}{0.52\textwidth}
\small
Almost every economic optimization is constrained: $c>0$, $n\in[0,1]$, $a'\ge-\underline a$. At an optimum the Karush--Kuhn--Tucker conditions hold:
$\nabla_x\mathcal{L}=0$, $\mu\ge0$, $h(x)\ge0$, and
\[
  \mu\,h(x)=0 .
\]
It is the last one --- \textbf{complementary slackness} --- that is awkward numerically: it is not an equation you can hand to a solver.
\medskip

The Fischer--Burmeister function turns it into one:
\[
  \phi(a,b)=a+b-\sqrt{a^{2}+b^{2}},
\]
\vspace{-3mm}
\[
  \phi(a,b)=0 \iff a,b\ge0,\; ab=0 .
\]
\end{column}
\begin{column}{0.45\textwidth}
\begin{conceptbox}[title=\textbf{Where you met this}]\small
2.A1 used exactly this $\phi$ to write the borrowing-constrained household's optimality conditions as a square system --- and pointed here for how to solve it.
\end{conceptbox}
\medskip
\begin{takeawaybox}\footnotesize
With $\phi$ in hand the KKT system becomes $F(x,\mu)=0$, and 4.A1's machinery applies directly: Newton, a Jacobian, a line search.\\[1mm]
$\phi$ is not differentiable at the origin --- exactly where the constraint is marginally binding --- so solvers smooth it and drive the smoothing to zero.
\end{takeawaybox}
\end{column}
\end{columns}
\vspace{1mm}
\src{$\phi(a,b)=a+b-\sqrt{a^2+b^2}$ in the instructor's own notation, from QM Lec.~8, ``Handling Inequality Constraints''.}
\end{frame}

%=============================================================================
\begin{frame}{Choosing, and a Warning About Worked Examples}
\begin{columns}[T]
\begin{column}{0.51\textwidth}
\renewcommand{\arraystretch}{1.15}
\begin{center}\footnotesize
\begin{tabularx}{\linewidth}{@{}>{\raggedright\arraybackslash}p{1.9cm} >{\raggedright\arraybackslash}X@{}}
\toprule
\textbf{Method} & \textbf{Use it when} \\
\midrule
Newton & $N$ small, $H$ cheap and reliable \\
BFGS & the default for smooth $f$ up to a few hundred parameters \\
L-BFGS & thousands of parameters; only $O(N)$ memory \\
Nelder--Mead & $f$ noisy, kinked, or a black box \\
\bottomrule
\end{tabularx}
\end{center}
\end{column}
\begin{column}{0.46\textwidth}
\begin{cautionbox}[title=\textbf{Check the arithmetic in your sources}]\footnotesize
Recomputing the conjugate-gradient example that accompanies this material turned up a printed direction
$d_1=[-80/81,\,\mathbf{40/81}]$ whose second component should be $\mathbf{20/81}$ --- and with the printed value the next two lines of the derivation do not follow.
\end{cautionbox}
\medskip
\begin{takeawaybox}\footnotesize
Worked examples are the most useful thing in a numerical-methods course and the easiest place for an error to survive, because nobody re-runs arithmetic that is already typeset. \textbf{Re-run it.}
\end{takeawaybox}
\end{column}
\end{columns}
\end{frame}

%=============================================================================
\begin{frame}{Takeaways}
\begin{takeawaybox}
\begin{itemize}\footnotesize
  \item In $\mathbb{R}^{N}$ there is no bracket to shrink. A \textbf{local quadratic model} replaces it, and the Hessian is what makes that expensive: a million entries and 333 million flops per step at $N=1{,}000$.
  \item The \textbf{secant condition} $Hs=y$ is free information from the step you just took. BFGS imposes it exactly --- verified to $4\times10^{-16}$, not ``approximately''.
  \item Direction and step length are separate problems; the step length is \textbf{4.A2, nested}.
  \item \textbf{Fast and accurate are different questions.} On Rosenbrock, BFGS used half the evaluations of Nelder--Mead and finished five orders of magnitude further away.
  \item Binding constraints are handled by turning complementary slackness into an equation --- Fischer--Burmeister --- after which everything reduces to 4.A1.
  \item When $N$ is enormous the Hessian stops fitting --- $10^{12}$ entries at $N=10^{6}$ --- so \textbf{CG} and \textbf{L-BFGS} replace it; when the objective is a sum over millions of observations, \textbf{SGD} replaces the gradient. Both return in Session 7.
  \item When the objective is \textbf{multimodal}, every method here stops at the first optimum it meets. Annealing escapes by accepting worse points on purpose --- and the same accept--reject rule, at fixed temperature, is how you \emph{sample} a posterior (Session 9).
\end{itemize}
\end{takeawaybox}
\end{frame}

%=============================================================================
\begin{frame}{Bridge: All of This Assumed You Had Derivatives}
\begin{columns}[T]
\begin{column}{0.53\textwidth}
\small
Newton needs a gradient and a Hessian. BFGS needs a gradient. The Jacobian in 4.A1 needs $N^{2}$ partial derivatives. Every method in the last two decks has quietly assumed that $\nabla f$ arrives from somewhere.
\medskip

For a model you wrote down, it does --- you can differentiate by hand, once, and check it. For a model with a hundred equations, or a solver nested inside another solver, hand-differentiation stops being realistic. And the obvious substitute, a finite difference, has a failure mode we have already measured.
\end{column}
\begin{column}{0.44\textwidth}
\begin{conceptbox}[title=\textbf{Next (4.A4)}]\small
Where derivatives come from: forward and central differences, the step size that is optimal and why, the complex-step trick that has no cancellation at all, and automatic differentiation.\\[1mm]
Then Newton, run three ways, with the Jacobian supplied by each.
\end{conceptbox}
\end{column}
\end{columns}
\end{frame}

\end{document}
